\chapter{Pressure Sores (Bedsores)}

\section*{Definition}
Pressure sores (bedsores, decubitus ulcers) are localized injuries to the skin and underlying tissue caused by prolonged pressure. They most commonly develop over bony prominences in immobile people. They are staged from I (non-blanchable erythema) to IV (full-thickness tissue loss).

\section*{What Happens in Your Body}
Sustained pressure exceeding capillary filling pressure (32 mmHg) compresses blood vessels, causing ischemia, hypoxia, and tissue necrosis. Shear and friction forces contribute by damaging superficial tissues. Moisture from incontinence or sweating macerates the skin, increasing vulnerability. Pressure sores are graded by depth of tissue involvement.

\section*{Causes}
\begin{itemize}
\item Prolonged immobility (bedridden, wheelchair-bound)
\item Neurological impairment (spinal cord injury, stroke, dementia)
\item Advanced age (thinner skin, decreased subcutaneous fat)
\item Malnutrition and dehydration
\item Incontinence (urinary or fecal)
\item Sensory loss (inability to feel pressure or pain)
\item Friction and shear forces from repositioning
\item Moisture from sweat, incontinence, or wound drainage
\item Comorbidities: diabetes, peripheral vascular disease, anemia
\end{itemize}

\section*{When to See a Doctor}
\begin{itemize}
\item Non-blanchable erythema of intact skin (Stage I)
\item Partial-thickness skin loss with blister or abrasion (Stage II)
\item Full-thickness skin loss with visible subcutaneous fat (Stage III)
\item Full-thickness tissue loss with exposed bone, muscle, or tendon (Stage IV)
\item Pain or tenderness over bony prominences
\item Location over sacrum, heels, elbows, hips, shoulders
\item Signs of infection: purulent drainage, foul odor, surrounding cellulitis
\end{itemize}

\section*{Related Symptoms}
\begin{itemize}
\item Cellulitis and soft tissue infection
\item Osteomyelitis (bone infection)
\item Sepsis
\item Wound infection
\item Malnutrition
\item Dehydration
\end{itemize}

\section*{Self-Care / Home Management}
\begin{itemize}
\item Reposition every 2 hours in bed or every 15 minutes in a wheelchair
\item Use pressure-relieving surfaces (specialty mattresses, cushions)
\item Keep skin clean and dry
\item Apply barrier creams for incontinence
\item Maintain adequate nutrition and hydration
\item Inspect skin daily for early signs of breakdown
\item Avoid massaging bony prominences
\end{itemize}

\section*{How Your Doctor Will Evaluate You}
\begin{itemize}
\item Clinical staging of the pressure sore
\item Wound culture if signs of infection
\item Complete blood count, albumin, prealbumin levels
\item X-ray or MRI if osteomyelitis suspected
\item Wound biopsy if malignancy suspected
\end{itemize}

\section*{Treatment Options}
\begin{itemize}
\item Pressure offloading (repositioning, specialty surfaces)
\item Wound debridement of necrotic tissue
\item Appropriate wound dressings (foam, hydrocolloid, alginate)
\item Negative pressure wound therapy for Stage III/IV
\item Antibiotics for cellulitis or osteomyelitis
\item Surgical closure (flap or graft) for severe wounds
\item Nutritional support (high-protein diet, supplements)
\end{itemize}

\section*{References}
\begin{thebibliography}{99}
\bibitem{pressure_sores:ref1} National Pressure Ulcer Advisory Panel, European Pressure Ulcer Advisory Panel, Pan Pacific Pressure Injury Alliance. \textit{Prevention and Treatment of Pressure Ulcers: Clinical Practice Guideline}. Cambridge Media, 2019.
\bibitem{pressure_sores:ref2} Edsberg LE, Black JM, Goldberg M, et al. ``Revised National Pressure Ulcer Advisory Panel pressure injury staging system.'' \textit{Journal of Wound, Ostomy and Continence Nursing}, 2016;43(6):585--597.
\bibitem{pressure_sores:ref3} Reddy M, Gill SS, Rochon PA. ``Preventing pressure ulcers: a systematic review.'' \textit{Journal of the American Medical Association}, 2006;296(8):974--984.
\bibitem{pressure_sores:ref4} Lyder CH, Ayello EA. ``Pressure ulcers: a patient safety issue.'' In: \textit{Patient Safety and Quality: An Evidence-Based Handbook for Nurses}. Agency for Healthcare Research and Quality, 2008.
\bibitem{pressure_sores:ref5} Chou R, Dana T, Bougatsos C, et al. ``Pressure ulcer risk assessment and prevention: a systematic comparative effectiveness review.'' \textit{Annals of Internal Medicine}, 2013;159(1):28--38.
\bibitem{pressure_sores:ref6} Smith ME, Totten A, Hickam DH, et al. ``Pressure ulcer treatment strategies: a systematic comparative effectiveness review.'' \textit{Annals of Internal Medicine}, 2013;159(1):39--50.

\end{thebibliography}
